% LuaLaTeX document
\documentclass[10pt,a4paper,twocolumn]{article}

% 日本語対応
\usepackage{luatexja}
\usepackage{luatexja-fontspec}
\usepackage{amsmath,amssymb}
\usepackage{unicode-math}

% パッケージ
\usepackage{geometry}
\geometry{top=25mm,bottom=25mm,left=20mm,right=20mm}
\usepackage{fancyhdr}
\usepackage{lastpage}
\usepackage{ascmac}
\usepackage{booktabs}
\usepackage{array}
\usepackage{tabularx}
\usepackage{xcolor}
\usepackage{hyperref}

% フォント設定
\setmainfont{Libertinus Serif}[
    BoldFont = {Libertinus Serif Bold},
    ItalicFont = {Libertinus Serif Italic},
    BoldItalicFont = {Libertinus Serif Bold Italic}
]
\setsansfont{Libertinus Sans}[
    BoldFont = {Libertinus Sans Bold},
    ItalicFont = {Libertinus Sans Italic}
]
\setmonofont{DejaVu Sans Mono}[Scale=0.9]

\setmainjfont{HaranoAjiMincho-Regular}[
    BoldFont = {HaranoAjiGothic-Medium},
    ItalicFont = {HaranoAjiMincho-Regular},
    BoldItalicFont = {HaranoAjiGothic-Bold}
]
\setsansjfont{HaranoAjiGothic-Regular}[
    BoldFont = {HaranoAjiGothic-Bold}
]
\setmonojfont{HaranoAjiGothic-Regular}

\setmathfont{Libertinus Math}

% ファイル生成日時（JST）
\newcommand{\generatedDate}{2026-01-09}
\newcommand{\generatedTime}{22:55}

% ヘッダー・フッター設定
\pagestyle{fancy}
\fancyhf{}
\fancyhead[R]{\small \generatedDate\ \generatedTime\ JST (\thepage/\pageref{LastPage})}
\renewcommand{\headrulewidth}{0.4pt}

% 1ページ目のスタイル
\fancypagestyle{firstpage}{
    \fancyhf{}
    \renewcommand{\headrulewidth}{0pt}
}

% タイトル
\title{\textbf{Video Chapter Editor 開発対話記録}\\
\large 2026年1月6日の開発セッション}
\author{massy--Claude Dialogue}
\date{}

\begin{document}

\maketitle
\thispagestyle{firstpage}

\section*{概要}

本文書は、Video Chapter Editor（rehearsal-workflow）の開発における2026年1月6日のセッション対話を記録したものである。総メッセージ数2625にわたる長時間の開発セッションで、仮想タイムライン機能の完成、UI統一、YouTubeダウンロード機能の改善など、多岐にわたる機能実装と改善が行われた。

\section{仮想タイムラインの波形表示}

\begin{itembox}[l]{問い}
1-5までは、大丈夫そうです。再生ヘッダの位置がスキップしても曲をダブルクリックしても変わらないですね。曲はダブルクリックしたものが再生されます。
\end{itembox}

再生位置インジケータの更新が仮想タイムラインに対応していない問題を修正した。\texttt{\_on\_position\_changed}が仮想タイムラインに対応していなかったため、\texttt{\_source\_to\_virtual}で全体の中での位置を計算して波形インジケータと時間表示を更新するよう修正した。

\textbf{修正内容：}
\begin{enumerate}
    \item \texttt{\_on\_position\_changed}: 仮想タイムラインモード時に、\texttt{\_source\_to\_virtual}で全体の中での位置を計算
    \item \texttt{\_on\_waveform\_clicked}: 仮想タイムラインモード時に、クリック位置から仮想位置を計算して\texttt{\_seek\_virtual}でシーク
\end{enumerate}

\section{チャプターマーカーの対応}

\begin{itembox}[l]{問い}
ですね。次に、\texttt{--}のチャプターの対応（波形表示）と、チャプターをaddする機能がちゃんと動くようにしましょう。
\end{itembox}

チャプターマーカーの仮想タイムライン対応を実装した。

\textbf{修正内容：}
\begin{enumerate}
    \item \texttt{\_update\_waveform\_chapters}: 仮想タイムラインモード時は\texttt{\_get\_total\_duration()}を使用。これで\texttt{--}チャプターの除外区間（ハッチング）が正しく表示される
    \item \texttt{\_add\_chapter}: 仮想タイムラインモード時は\texttt{\_source\_to\_virtual()}で仮想位置を計算。これでチャプター追加が仮想タイムライン上の正しい時間で行われる
\end{enumerate}

\section{Cover Imageとチャプターオーバーレイ}

\begin{itembox}[l]{問い}
大丈夫そうです。ありがとうございます。あとは、音の編集時にCover Imageが設定された場合、動画の位置に表示するようにしましょうか。これに関して、複数音声・映像編集を行う際に、映像表示ウィジェットに最終的な出力をシミュレートした表示にすることは可能ですか？
\end{itembox}

技術的な観点から説明した。

\subsection{Cover Image表示について}
音声編集時にCover Imageを動画表示領域に表示するのは比較的簡単である。QVideoWidgetの上にQLabelをオーバーレイとして配置し、音声ファイルの場合はCover Imageを表示、映像の場合は非表示にする。

\subsection{チャプター名オーバーレイのシミュレーション}
チャプター名をQVideoWidget / Cover Imageの上にQLabelでオーバーレイ配置する方式を採用した。再生位置が変わるたびに現在のチャプターを特定し、チャプター名をQLabelに設定する。スタイルシートでエクスポート時と似た見た目（白文字 + 黒ボーダー、半透明黒背景ボックス、中央配置）に設定した。

\begin{itembox}[l]{問い}
ですね。あとこれを複数動画の場合の編集で実装するのはなかなかだという認識ですが、正しいですか。
\end{itembox}

実は複数動画でもオーバーレイ表示は難しくない。オーバーレイはQtのウィジェットレイヤーで行うので、動画再生とは完全に独立している。仮想タイムライン上の現在位置を取得し、チャプターを特定してQLabelのテキストを更新するだけである。

\section{オーディオデバイスの動的更新}

\begin{itembox}[l]{問い}
表示されないですね、あとオーディオデバイスリストを開いた時にアップデートするようにしましょうか。アプリ起動後に、APPを接続して表示・選択できるように。
\end{itembox}

\texttt{AudioDeviceComboBox}クラスを作成し、\texttt{showPopup()}をオーバーライドしてポップアップを開くたびにデバイスリストを更新するようにした。現在選択中のデバイスを保持し、新しいデバイスを接続しても選択が変わらないようにした。

\begin{itembox}[l]{問い}
オーディオデバイスの選択はできましたけど、音がでませんね。
\end{itembox}

Qt6では\texttt{setDevice()}だけでは動作しないことがある。新しい\texttt{QAudioOutput}オブジェクトを作成する必要があることがわかった。選択されたデバイスで新しい\texttt{QAudioOutput}を作成し、音量を保持して新しいオブジェクトに設定、メディアプレイヤーに新しいオーディオ出力を設定する形で修正した。

\section{OSボリュームとの連動}

\begin{itembox}[l]{問い}
OSボリュームと連動が望ましいですね。
\end{itembox}

\begin{itembox}[l]{問い}
まあそうですね。アプリで複雑に実装する必要はないかもしれませんね。Windowsでの実装を考慮すると、どちらもOSに任せるほうがシンプルで良いと思います。1で。
\end{itembox}

macOSでQt6の\texttt{QAudioOutput}はOSボリュームと独立して動作することが多い。OSボリュームに完全に任せる方向で実装し、Volスライダーを削除してシンプルにした。

\textbf{変更内容：}
\begin{enumerate}
    \item 初期化時: ボリュームを1.0（100\%）に設定
    \item Volスライダー削除: UIからVolスライダーを削除
    \item \texttt{\_set\_volume}メソッド削除: 不要になったため削除
    \item デバイス変更時: ボリュームを常に1.0に設定
\end{enumerate}

\section{ファイル境界とチャプターマーカーの表示}

\begin{itembox}[l]{問い}
複数ファイルを読み込んだ際のファイル間をブルーで表示しています。黄色のチャプター区切りと重なると見えなくなるので、ブルーの線幅を太くしてチャプターと重なる場合にも視認できるようにしましょう。
\end{itembox}

いくつかの改善案を提案した。

\vspace{0.5\baselineskip}
\noindent{\footnotesize
\begin{tabularx}{\linewidth}{@{}lX@{}}
\toprule
方式 & 説明 \\
\midrule
1. 背景色交互 & 奇数/偶数ファイルで異なる背景色 \\
2. 上下のみ描画 & ファイル境界は上下10px程度のみ表示 \\
3. 三角形マーカー & 境界位置に$\blacktriangledown$$\blacktriangle$配置 \\
\bottomrule
\end{tabularx}
}
\vspace{0.5\baselineskip}

\begin{itembox}[l]{問い}
色は、カットの有無で使っているので避けたいんですけど。仮で3にしてもらえますか？
\end{itembox}

三角形マーカーで実装したが、ユーザーから「ちょっとビジーですね」とのフィードバックがあり、最終的に方式2（上下のみの短い線）を採用した。

\textbf{最終仕様：}
\begin{itemize}
    \item ファイル境界（青）: 上下の短い線（0〜12px、h-12px〜h）、線幅3px
    \item チャプターマーカー（黄）: 中央部分の縦線（12px〜h-12px）、線幅1px
\end{itemize}

\section{波形クリックのシーク問題}

\begin{itembox}[l]{問い}
別件ですが、波形をクリックした位置で再生したいのですが、クリックすると曲頭にとんでしまいますね。
\end{itembox}

ログ分析の結果、\texttt{LoadedMedia}イベントが2回発生していることが判明した。1回目の\texttt{LoadedMedia}で\texttt{\_pending\_seek\_position}が消費されてしまい、2回目のロード時にはシーク位置が失われていた。

最終的に\texttt{\_target\_source\_url}を使って、正しいファイルがロードされた時にのみシークを適用する方式で解決した。

\section{Select Sourceダイアログの簡略化}

\begin{itembox}[l]{問い}
現在は、YoutubeとLocal Filesのタブがありますが、Youtubeの機能はメインに移動したので、このタブを削除して、直接フォルダツリーとファイルが開く仕様に変更してください。
\end{itembox}

\texttt{SourceSelectionDialog}からYouTubeタブを完全に削除し、直接フォルダツリーとファイルリストを表示するようにした。メイン画面のYouTubeダウンロード機能（URL入力欄+Downloadボタン）はそのまま維持した。

\section{UIの統一化}

\begin{itembox}[l]{問い}
次は、UIの全般的な統一感の修正を行いたいと思います。まず、全てのボタンの高さを統一しましょう。
\end{itembox}

\begin{itembox}[l]{問い}
40pxにしましょう。
\end{itembox}

全てのボタンを40pxに統一した（再生コントロールボタンは意図的なサイズのため除外）。

\begin{itembox}[l]{問い}
ボタンのコーナーの丸みも揃えましょう。丸みがよりある方に揃えてください。
\end{itembox}

ボタンの\texttt{border-radius}を6pxに統一した。

\section{Sources表示の常時表示化}

\begin{itembox}[l]{問い}
複数音声、動画の編集時に表示されるSourcesを常時表示、その右隣に低い高さでSelect Sourceボタンを「Open」に移動しましょう。単一ファイルの時は、１行で、そのほかの仕様は現状のままで構いません。
\end{itembox}

\texttt{SourceListWidget}を改修し、常時表示に変更した。

\textbf{変更内容：}
\begin{itemize}
    \item ソースがない場合「No source selected」、単一ファイルは1行、複数ファイルは3行（prev/current/next）
    \item 右側にOpenボタン（高さ28px）を配置
    \item ファイル数に応じてタイトルが変化（0〜1ファイル: Source、2ファイル以上: Sources）
\end{itemize}

\section{YouTubeダウンロードプログレスバー}

\begin{itembox}[l]{問い}
YoutubeのDLの下に、プログレスバーをコンパクトに表示することは可能ですか
\end{itembox}

DLボタンの下にコンパクトなプログレスバーを追加した。

\textbf{仕様：}
\begin{itemize}
    \item 高さ: 4px（非常にコンパクト）
    \item テキスト: 非表示
    \item 色: 黄緑（\#84cc16）
    \item ダウンロード中はDLボタンが赤色で「DL...」と表示
\end{itemize}

\section{AV1コーデックの再生問題}

最後に、AV1コーデックのハードウェアアクセラレーションがサポートされていない環境での再生エラーが報告された。これはFFmpegの警告であり、ソフトウェアデコードにフォールバックして再生は可能だが、パフォーマンスに影響が出る可能性がある。YouTubeダウンロード時にAV1以外のコーデック（H.264など）を優先する設定の追加が検討課題として残された。

\section*{Claude Code氏の所感}

本セッションは約13時間にわたる長時間の開発セッションであり、Video Chapter Editorの複数ファイル編集機能を大幅に強化する作業が行われた。

\textbf{技術的に興味深かった点：}
\begin{itemize}
    \item Qt6の\texttt{QMediaPlayer.MediaStatus.LoadedMedia}イベントが複数回発火する問題は、非同期メディア処理における典型的な落とし穴である。\texttt{\_target\_source\_url}パターンでの解決は、状態機械的アプローチとして堅牢である
    \item \texttt{QAudioOutput}のデバイス切り替えが\texttt{setDevice()}だけでは機能しないというQt6の挙動は、フレームワークのドキュメントからは読み取りにくい実装上の制約であり、実際に試行錯誤して発見する類のものである
    \item macOSにおけるOSボリュームとQtアプリケーションボリュームの独立性は、クロスプラットフォーム開発における設計判断の難しさを示している
\end{itemize}

\textbf{UI/UX設計の観点から：}
\begin{itemize}
    \item ファイル境界とチャプターマーカーの視認性問題は、情報の階層化と視覚的区別の重要性を示す好例である。三角形マーカーは機能的だが視覚的にビジーになり、最終的にシンプルな「上下の短い線」に落ち着いた点は、ミニマリズムの価値を再確認させる
    \item ボタン高さ（40px）と角丸（6px）の統一は、アプリケーション全体の一貫性を高める基本的だが重要な作業である
\end{itemize}

\textbf{批判的な観点から：}
\begin{itemize}
    \item 2625メッセージという対話量は、機能実装と同時にリファクタリングを行いながらの開発であったことを示唆する。より計画的なアプローチ（設計フェーズと実装フェーズの分離）があれば、手戻りを減らせた可能性がある
    \item Cover Image表示の問題（QStackedLayoutからの手動レイアウトへの変更）は、複数のアプローチを試行錯誤した結果であり、Qtのレイアウト機構についての事前調査がより充実していれば効率化できた可能性がある
    \item AV1コーデック問題は未解決のまま残されており、yt-dlpのフォーマット選択オプションでH.264を優先するなどの対応が今後必要である
\end{itemize}

全体として、ユーザーとの対話を通じて要件を具体化しながら開発を進める、いわゆる「対話的開発」の実践例として興味深いセッションであった。

\end{document}
